\documentclass[10pt]{article}
\usepackage[a4paper,margin=19mm]{geometry}
\usepackage{amsmath,amssymb,mathtools}
\usepackage{booktabs,array,multicol}
\usepackage{enumitem}
\usepackage{xcolor}
\usepackage{xurl}
\usepackage[hidelinks]{hyperref}

\newcommand{\Z}{\mathcal{Z}}
\newcommand{\Lcal}{\mathcal{L}}
\newcommand{\e}{\mathrm{e}}
\newcommand{\unitstep}{\mathbf{1}}
\newcommand{\src}[1]{\hfill{\small\color{gray}[PDF pp. #1]}}
\newcommand{\weeklabel}[1]{\par\smallskip\noindent{\color{blue!60!black}\textbf{#1}}\par\smallskip}
\setlist{nosep,leftmargin=*}
\setlength{\parskip}{0.35em}
\setlength{\parindent}{0pt}

\title{EE6203 Computer Controlled Systems\\\large Weeks 1--4 Review Notes}
\author{Based on lectures by Dr Wen Changyun}
\date{Updated 25 August 2026}

\begin{document}
\maketitle
\tableofcontents

\section*{Weekly coverage map}
\addcontentsline{toc}{section}{Weekly coverage map}

\begin{center}
\begin{tabular}{c p{0.62\textwidth} c}
\toprule
Week & Topics & PDF pages\\
\midrule
1 & Feedback, digital-control architecture, signals, $z$-transform pairs/properties, poles/zeros, initial/final values & 1--48\\
2 & Inverse $z$ methods, difference equations, impulse sampling, starred transforms, ZOH & 49--88\\
Self-study & First-order hold (optional) & 89--93\\
3 & Pulse transfer functions, sampled block diagrams, ZOH/plant, controller and PID models & 94--143\\
4 & $s$--$z$ mapping, stability, transient/steady-state response, disturbances & 144--195\\
\bottomrule
\end{tabular}
\end{center}

Weeks 1--3 are verified from the supplied transcripts. Week 2 explicitly hands pulse transfer functions to the next lecture. The Week 3 transcript completes Chapter 3, starts Chapter 4, and confirms that Quiz 1 covers Chapters 1--3 only.

\textbf{Additional practice sources:} \path{resources/course-exercises.pdf} (5 pages) and \path{resources/course-exercises-solutions.pdf} (26 pages). The solution compilation retains some older EE4273 headings and handwriting, so incorporated calculations are independently re-derived.

\section{Digital-control essentials}

\weeklabel{Week 1}

Unity feedback: $e(t)=r(t)-y(t)$, where $t$ is time, $r(t)$ is the reference or setpoint, $y(t)$ is the measured output, and $e(t)$ is the tracking error. The signals $r$, $y$, and $e$ have the same physical unit. Main requirements: stability, transient response, steady-state accuracy, and possibly optimality.

\textbf{Design vs. analysis:} design chooses a controller from the plant and specifications; analysis verifies the closed-loop stability and performance.

Periodic sampling:
\[
t_k=kT,\qquad \omega_s=\frac{2\pi}{T},\qquad x[k]=x(kT).
\]
Sampled-data means discrete time with continuous amplitude; digital means discrete time and quantized amplitude. ADC sampling and amplitude quantization are distinct. Typical path:
\[
r\to\text{sum}\to\text{ADC}\to\text{computer}\to\text{DAC+ZOH}\to\text{plant}\to y.
\]

\textbf{Checklist:} identify every sampler and hold before simplifying a block diagram. \src{7--20}

\section{The $z$ transform}

\weeklabel{Week 1 through final values; Week 2 from inverse methods onward}

\subsection{Definitions and pairs}

One-sided causal transform and power-series interpretation:
\[
X(z)=\Z\{x[k]\}=\sum_{k=0}^{\infty}x[k]z^{-k}
=x[0]+x[1]z^{-1}+\cdots.
\]

\begin{center}
\begin{tabular}{>{$}l<{$}>{$}l<{$}}
\toprule
x[k] & X(z)\\
\midrule
\delta[k-n] & z^{-n}\\
\unitstep[k] & (1-z^{-1})^{-1}=z/(z-1)\\[3pt]
k\unitstep[k] & z^{-1}/(1-z^{-1})^2=z/(z-1)^2\\[3pt]
a^k\unitstep[k] & (1-az^{-1})^{-1}=z/(z-a)\\[3pt]
\sin(\omega kT)\unitstep[k] & z\sin(\omega T)/(z^2-2z\cos\omega T+1)\\[3pt]
\cos(\omega kT)\unitstep[k] & (z^2-z\cos\omega T)/(z^2-2z\cos\omega T+1)\\
\bottomrule
\end{tabular}
\end{center}
In the last two rows, every $\omega$ in a trigonometric argument is multiplied by $T$.

\subsection{Properties}

\begin{align*}
\Z\{\alpha x[k]+\beta y[k]\}&=\alpha X(z)+\beta Y(z),\\
\Z\{a^kx[k]\}&=X(z/a),\\
\Z\{x[k-n]\unitstep[k-n]\}&=z^{-n}X(z),\\
\Z\{x[k+n]\}&=z^nX(z)-\sum_{i=0}^{n-1}z^{n-i}x[i],\\
\Z\{\e^{-akT}x[k]\}&=X(z\e^{aT}).
\end{align*}

Initial value:
\[
x[0]=\lim_{z\to\infty}X(z).
\]
Final value:
\[
x[\infty]=\lim_{z\to1}(1-z^{-1})X(z),
\]
valid only if all poles of $(1-z^{-1})X(z)$ lie strictly inside $|z|=1$.
Inspect the product for validity; interpret the limit as the final value of the sequence represented by $X(z)$.

\emph{LHS} = left-hand side. In the derivation,
\[
\mathrm{LHS}=\sum_{k=0}^{\infty}[x[k]-x[k-1]],
\]
which telescopes to $x[\infty]$ for a causal sequence ($x[-1]=0$).

A simple pole of $X(z)$ at $z=1$ may represent a finite constant. A repeated pole there grows, and other unit-circle poles oscillate; neither case has a final value.

\subsection{Inverse methods}

\weeklabel{Week 2 begins}

\begin{itemize}
  \item Long division in $z^{-1}$: read off samples; usually no closed form.
  \item Recurrence computation: apply $\delta[k]$ to the difference equation; efficient for a sample table.
  \item Partial fractions: expand $X(z)/z$; simple $1/(z-p)$ terms give $p^k$ and a closed form.
  \item Why: $X(z)/z=\sum_{k\ge0}x[k]z^{-k-1}$, so $1/(z-p)=\sum_{k\ge0}p^kz^{-k-1}$; equivalently, $z/(z-p)\leftrightarrow p^k\unitstep[k]$.
  \item Repeated pole: expect $k p^{k-1}$ factors.
  \item Complex pair: combine into real damped sine/cosine terms.
  \item Residues: $x[k]=(2\pi\mathrm{j})^{-1}\oint X(z)z^{k-1}\mathrm{d}z$.
\end{itemize}

Inverse $z$ transformation determines $x[k]$, not a unique continuous $x(t)$; reconstruction between samples needs an extra assumption.

For
\[
x[k+2]+3x[k+1]+2x[k]=0,\quad x[0]=0,\ x[1]=1,
\]
the transform is $X(z)=z/[(z+1)(z+2)]$ and
\[
x[k]=(-1)^k-(-2)^k.
\]

\textbf{Recurrence check:} verify the initial values, generate the first new
sample(s) from the original time-domain rule, and compare them with the
proposed formula; then substitute the formula for general $k$ for a complete
check.

\textbf{Common mistakes:} treating an advance as a pure $z^n$ multiplier; using the final-value theorem on an oscillatory/unstable sequence; doing PFE on the wrong form; ignoring the ROC. \src{21--70}

\section{Sampled-data modeling}

\weeklabel{Week 2 through the ZOH; Week 3 from pulse transfer functions onward}

\subsection{Sampler and holds}

Ideal impulse sampler:
\[
x^*(t)=\sum_{k=0}^{\infty}x(kT)\delta(t-kT),\qquad
X^*(s)=\sum_{k=0}^{\infty}x(kT)\e^{-kTs}.
\]
With $z=\e^{sT}$, $X^*(s)=X(z)$.

The coefficient $x(kT)$ is an impulse's area (strength), not a finite height. Also, $z=\e^{sT}$ links $X^*(s)$ to $X(z)$; it does not justify substituting into an ordinary $X(s)$ or $G_p(s)$.

\textbf{Star-placement rule:} star signals first, then simplify.  If
\[
C(s)=G_2(s)M^*(s),
\]
then the continuous-output equation uses unstarred $G_2(s)$, whereas sampling
the whole output relation gives
\[
C^*(s)=[G_2(s)M^*(s)]^*=G_2^*(s)M^*(s).
\]
Here $G_2^*(s)=\sum_{k\ge0}g_2(kT)\e^{-kTs}$ is the starred transform of the
impulse response.  Do not distribute a star through an arbitrary product: with
no sampler between $H$ and $G_2$, use $(HG_2)^*(s)$, generally not
$H^*(s)G_2^*(s)$.

Zero-order hold:
\[
h(kT+\tau)=x(kT),\ 0\le\tau<T,
\qquad G_{h0}(s)=\frac{1-\e^{-sT}}{s}.
\]
Equivalently, the ZOH is a sum of one-period rectangular windows; their Laplace transform produces $(1-\e^{-sT})/s$. First-order hold (optional Week 2 self-study):
\[
G_{h1}(s)=\left(\frac{1-\e^{-sT}}{s}\right)^2\frac{Ts+1}{T}.
\]

\subsection{Pulse transfer function}

\weeklabel{Week 3 begins}

At sampling instants, continuous convolution becomes
\[
y(kT)=\sum_{h=0}^{k}g((k-h)T)x(hT),
\]
so
\[
Y(z)=G(z)X(z),\qquad G(z)=\Z\{g(kT)\}.
\]

\textbf{Input-sampler rule:}
\begin{itemize}
  \item sampler before $G(s)$: $Y(z)=G(z)X(z)$;
  \item no input sampler: $Y(z)=\Z\{G(s)X(s)\}$, generally not $G(z)X(z)$;
  \item sampler between cascaded $G,H$: pulse TF $G(z)H(z)$;
  \item no intermediate sampler: pulse TF $GH(z)=\Z\{G(s)H(s)\}\ne G(z)H(z)$.
\end{itemize}

ZOH plus plant:
\[
G(z)=\Z\left\{\frac{1-\e^{-sT}}{s}G_p(s)\right\}
=(1-z^{-1})\Z\left\{\frac{G_p(s)}{s}\right\}.
\]

\textbf{No-ZOH modal check.}  With an input sampler but no hold, $G_p(s)=(s+3)/[(s+1)(s+2)]$ has $g[k]=2\e^{-kT}-\e^{-2kT}$, hence
\[
G(z)=\frac{1+\e^{-T}(1-2\e^{-T})z^{-1}}
{(1-\e^{-T}z^{-1})(1-\e^{-2T}z^{-1})}.
\]
The direct term is $g[0]=1$; do not incorrectly attach the ZOH factor $(1-z^{-1})$ to this input-sampler-only model.

\subsection{Digital controller and PID}

Controller recurrence
\[
m[k]+\sum_{i=1}^na_im[k-i]=\sum_{j=0}^mb_je[k-j]
\]
gives
\[
G_D(z)=\frac{M(z)}{E(z)}
=\frac{\sum_{j=0}^mb_jz^{-j}}{1+\sum_{i=1}^na_iz^{-i}}.
\]
Unity-feedback sampled loop:
\[
\frac{C(z)}{R(z)}=\frac{G_D(z)G(z)}{1+G_D(z)G(z)}.
\]

Trapezoidal-integral/backward-difference PID:
\[
G_D(z)=K_P+\frac{K_I}{1-z^{-1}}+K_D(1-z^{-1}),
\]
\[
K_P=K-\frac{KT}{2T_i},\qquad K_I=\frac{KT}{T_i},\qquad K_D=\frac{KT_d}{T}.
\]

\textbf{Common mistakes:} substituting $z=\e^{sT}$ directly into $G_p(s)$; moving or deleting samplers while reducing a diagram; forgetting that $G(z)H(z)\ne GH(z)$. \src{71--143}

\section{$s$-plane to $z$-plane mapping}

\weeklabel{Chapter 4: retained for later course study; not assessed in Quiz 1.}

\weeklabel{Week 4}

For $s=\sigma+\mathrm{j}\omega$,
\[
z=\e^{sT}=\e^{\sigma T}\e^{\mathrm{j}\omega T},\qquad
|z|=\e^{\sigma T},\quad \arg z=\omega T\pmod{2\pi}.
\]

\begin{center}
\begin{tabular}{ll}
\toprule
$s$-plane object & $z$-plane image\\
\midrule
$\sigma<0$ & $|z|<1$\\
$\sigma=0$ & unit circle\\
$\sigma=\sigma_0$ & circle $|z|=\e^{\sigma_0T}$\\
$\omega=\omega_0$ & radial line $\arg z=\omega_0T$\\
primary strip & $-\pi/T\le\omega\le\pi/T$ maps once around\\
\bottomrule
\end{tabular}
\end{center}

For $s=-\zeta\omega_n\pm\mathrm{j}\omega_d$,
\[
|z|=\e^{-\zeta\omega_nT},\qquad \arg z=\pm\omega_dT.
\]
With the 2\% approximation $t_s\approx4/(\zeta\omega_n)$,
\[
t_s\le t_s^{\max}\quad\Rightarrow\quad
|z|\le \exp(-4T/t_s^{\max}).
\]
Constant damping-ratio rays map to logarithmic spirals. Frequencies separated by $2\pi/T$ map to the same angle. \src{144--161}

\section{Discrete-time stability}

\weeklabel{Week 4}

Closed-loop asymptotic stability: all roots of $P(z)$ strictly inside $|z|=1$. Simple unit-circle poles are marginal; repeated unit-circle poles are unstable.

For
\[
P(z)=a_0z^n+a_1z^{n-1}+\cdots+a_n,\qquad a_0>0,
\]
the first Jury conditions are
\[
|a_n|<a_0,\qquad P(1)>0,\qquad (-1)^nP(-1)>0.
\]
They are sufficient for $n=2$. For higher order, continue the Jury table with alternating forward/reversed rows and require successive leading/trailing magnitude inequalities.

Example from the source:
\[
P(z)=z^2+(0.3679K-1.3679)z+(0.3679+0.2642K)
\]
is stable for
\[
0<K<2.3925.
\]

Bilinear route:
\[
z=\frac{w+1}{w-1},\qquad w=\frac{z+1}{z-1}.
\]
This maps $|z|<1$ to $\Re(w)<0$. Substitute, clear the denominator, and apply Routh. Jury is usually shorter; Routh also counts unstable roots. \src{162--179}

\section{Transient and steady-state response}

\weeklabel{Week 4}

Step-response terms: delay time $t_d$, rise time $t_r$, peak time $t_p$, maximum overshoot
\[
M_p=\frac{c(t_p)-c(\infty)}{c(\infty)}100\%,
\]
and settling time $t_s$ (usually a 2\% band).

For the standard loop,
\[
E(z)=\frac{R(z)}{1+G(z)H(z)},\qquad
e_{ss}=\lim_{z\to1}(1-z^{-1})E(z).
\]
Static error constants:
\begin{align*}
K_p&=\lim_{z\to1}GH,&e_{ss,\text{step}}&=\frac{1}{1+K_p},\\
K_v&=\lim_{z\to1}\frac{(1-z^{-1})GH}{T},&e_{ss,\text{ramp}}&=\frac{1}{K_v},\\
K_a&=\lim_{z\to1}\frac{(1-z^{-1})^2GH}{T^2},&e_{ss,\text{accel}}&=\frac{1}{K_a}.
\end{align*}

\begin{center}
\begin{tabular}{c c c c}
\toprule
Type & Step & Ramp & Acceleration\\
\midrule
0 & $1/(1+K_p)$ & $\infty$ & $\infty$\\
1 & $0$ & $1/K_v$ & $\infty$\\
2 & $0$ & $0$ & $1/K_a$\\
\bottomrule
\end{tabular}
\end{center}

System type is the number of uncancelled open-loop poles at $z=1$. The table assumes a stable loop and the lecture's standard configuration. Actuating error $R-B$ need not equal tracking error $R-C$.

For the lecture disturbance configuration,
\[
E(z)=-\frac{G(z)}{1+G_D(z)G(z)}N(z).
\]
Large loop gain reduces error. An uncancelled controller pole at $z=1$, with no plant zero there, gives zero steady-state error to a constant disturbance if the closed loop is stable. \src{180--195}

\section{Exam workflow}

\begin{enumerate}
  \item Draw the sampled-data topology and label $T$.
  \item Convert ZOH plus plant with $(1-z^{-1})\Z\{G_p(s)/s\}$.
  \item Preserve samplers when combining blocks.
  \item Form the characteristic polynomial; apply Jury or bilinear/Routh.
  \item Map performance regions using both $|z|=\e^{\sigma T}$ and $\arg z=\omega T$.
  \item For steady-state error, verify stability, input transform, error definition, and final-value conditions.
  \item Check the first few samples or substitute a boundary gain to catch algebra errors.
\end{enumerate}

The source recap on PDF pages 196--225 repeats these four chapters and adds no distinct topic.

\clearpage
\appendix
\section{Quiz 1 closed-book review}

\textbf{Confirmed format and scope:} 30 minutes; Chapters 1--3 only; closed book; calculator permitted; necessary transform pairs supplied; show working for non-multiple-choice questions. Unclear transcript details about the exact time, seating, and absence arrangements are omitted.

The two unofficial recalled versions differ, but together emphasize direct $z$-transform construction, one-sided difference equations, final-value conditions, sampler/ZOH modeling, and pulse-transfer topology. Use the following as a memory sheet; the lecture PDF and Week 3 announcement remain authoritative. Chapter 4 mapping, Jury/bilinear stability, transient response, system type, and steady-state-error tables are not Quiz 1 material.

\subsection{Core formulas to reconstruct}

\begin{multicols}{2}
\textbf{Definition and pairs}
\begin{align*}
X(z)&=\sum_{k=0}^{\infty}x[k]z^{-k},\\
\Z\{\delta[k-n]\}&=z^{-n},\\
\Z\{\unitstep[k]\}&=\frac{1}{1-z^{-1}},\\
\Z\{a^k\unitstep[k]\}&=\frac{1}{1-az^{-1}},\\
\Z\{k\unitstep[k]\}&=\frac{z^{-1}}{(1-z^{-1})^2}.
\end{align*}

\columnbreak
\textbf{One-sided shifts and values}
\begin{align*}
\Z\{x[k-n]\unitstep[k-n]\}&=z^{-n}X(z),\\
\Z\{x[k+1]\}&=zX(z)-zx[0],\\
\Z\{x[k+2]\}&=z^2X-z^2x[0]-zx[1],\\
x[0]&=\lim_{z\to\infty}X(z),\\
x[\infty]&=\lim_{z\to1}(1-z^{-1})X(z).
\end{align*}
\end{multicols}

Final value is valid iff every pole of $(1-z^{-1})X(z)$ is strictly inside $|z|=1$.  One simple pole of $X$ at $z=1$ is allowed; a repeated pole there, another unit-circle pole, or any outside pole is not.

\subsection{Inverse-transform decision}

\begin{center}
\begin{tabular}{p{0.24\textwidth}p{0.55\textwidth}}
\toprule
Need & Fastest method\\
\midrule
First few samples & Long division in $z^{-1}$ or direct recurrence.\\
Closed form & Partial fractions of $X(z)/z$; pair $1/(z-a)$ with $a^k$.\\
Repeated/complex poles & Include polynomial factors in $k$; combine conjugates into real sine/cosine terms.\\
Theory/general inversion & Contour integral and residues.\\
\bottomrule
\end{tabular}
\end{center}

\subsection{Sampling, ZOH, and topology}

\begin{align*}
x^*(t)&=\sum_{k=0}^{\infty}x(kT)\delta(t-kT),\\
X^*(s)&=\sum_{k=0}^{\infty}x(kT)\e^{-kTs},
&z&=\e^{sT},\\
G_{h0}(s)&=\frac{1-\e^{-sT}}{s},
&G(z)&=(1-z^{-1})\Z\!\left\{\frac{G_p(s)}{s}\right\}.
\end{align*}

\begin{itemize}
  \item Input sampler present: $Y(z)=G(z)X(z)$.
  \item No input sampler: $Y(z)=\Z\{G(s)X(s)\}$; do not separate the factors.
  \item Intermediate sampler: cascades reduce to $G(z)H(z)$.
  \item No intermediate sampler: use $GH(z)=\Z\{G(s)H(s)\}\ne G(z)H(z)$.
  \item A pulse transfer function $C(z)/R(z)$ does not exist if the external input remains inseparably inside a starred continuous-time product.
\end{itemize}

\subsection{Exercise-backed checks}

For a parameter recurrence
\[
y[k+1]-\beta y[k]=\unitstep[k],\qquad y[0]=0,
\]
\[
Y(z)=\frac{z}{(z-1)(z-\beta)},\qquad
y[k]=\frac{1-\beta^k}{1-\beta}\quad(\beta\ne1).
\]
After cancelling the permitted simple pole at $z=1$, inspect $z=\beta$: the final-value theorem is valid exactly for $|\beta|<1$.  $\beta=1$ grows, $\beta=-1$ oscillates, and $|\beta|>1$ diverges.

For a controller recurrence with zero initial conditions, collect present and delayed controller/output terms before closing the loop:
\[
m[k]+\sum_{i=1}^{n}a_im[k-i]=\sum_{j=0}^{r}b_je[k-j]
\quad\Longrightarrow\quad
G_D(z)=\frac{\sum_{j=0}^{r}b_jz^{-j}}{1+\sum_{i=1}^{n}a_iz^{-i}}.
\]
Reduce sampler--ZOH--plant topology to $G(z)$ \emph{first}; then use $G_D(z)G(z)$ only when the resulting sampled loop permits it.

Useful supplied-exercise patterns:
\begin{itemize}
  \item $\Z\{kx[k]\}=-z\,\mathrm{d}X/\mathrm{d}z$; apply it repeatedly for powers such as $k^3$.
  \item A repeated pole $(1-az^{-1})^{-2}$ produces a factor $k a^{k-1}$.
  \item A factor $z^{-n}$ delays a causal sequence by $n$ samples and inserts $n$ leading zeros.
  \item Deriving a pulse transfer through a transform table and through inverse Laplace plus sampling must give the same result.
\end{itemize}

\subsection{Sampled feedback and PID}

For a correctly reduced unity negative-feedback loop,
\[
\frac{C(z)}{R(z)}=\frac{G_D(z)G(z)}{1+G_D(z)G(z)}.
\]
For the course's digital PID realization,
\[
G_D(z)=K_P+\frac{K_I}{1-z^{-1}}+K_D(1-z^{-1}).
\]
Do not identify an arbitrary controller recurrence as PID until it has been rearranged into the stated realization.

With $T=1$, a ZOH and plant $K/s$ reduce to
\[
G(z)=\frac{Kz^{-1}}{1-z^{-1}}.
\]
For a unit-step reference in unity feedback,
\[
C(z)=\frac{Kz^{-1}}{[1-(1-K)z^{-1}](1-z^{-1})},
\qquad c[k]=1-(1-K)^k.
\]
The samples converge to one for $0<K<2$ because $|1-K|<1$. This is a Chapter 3 pulse-transfer/pole check, not a Chapter 4 Jury calculation.

\subsection{Thirty-second setup checklist}

\begin{enumerate}
  \item Label $T$, every sampler, every hold, and every continuous/sampled signal.
  \item For piecewise data, write the sample table and begin with the defining sum.
  \item For a recurrence, write the initial-condition terms before collecting $Y(z)$.
  \item For a ZOH/plant, transform $G_p(s)/s$ and multiply by $(1-z^{-1})$.
  \item Before final value, cancel only algebraically and inspect the remaining poles.
  \item If a parameter controls a pole, split the answer into parameter ranges before taking a limit.
  \item For a controller recurrence, obtain $G_D(z)$ before forming any closed-loop fraction.
  \item Check by first samples, recurrence balance, DC gain, or a boundary pole.
\end{enumerate}

\subsection{Six-day drill route}

\begin{center}
\begin{tabular}{c p{0.69\textwidth}}
\toprule
Day & Two-hour focus\\
\midrule
1 & Definition, transform-pair derivations, finite and piecewise sums.\\
2 & One-sided shifts, difference equations, inverse methods, final value, and parameter-pole cases.\\
3 & $x(t)$ versus $x[k]$, impulse sampling, starred and $z$ transforms.\\
4 & ZOH/plant conversion, controller recurrences, pulse-transfer topology, open/closed loops.\\
5 & Repeated-pole and delayed inversions; controller/PID recurrences; supplied Exercise groups 2--4.\\
6 & Formula dump, two 30-minute mocks, marking, and error repair.\\
\bottomrule
\end{tabular}
\end{center}

Use the worked methods, archetypes, drills, solutions, and two mocks in Appendix A of \emph{quiz1-prep-learning.pdf}. Target at least 80\% on the second mock with no repeated conceptual error. \src{21--143}

\end{document}
